\section*{第 \ref{ch:g3:sun-days-seasons} 章 --- 天空中的太阳：昼夜与季节}

\begin{solution}{exo:g3:sun-days-seasons:1}
夏季时太阳最高、路径最长，白天很长；
冬季时太阳最低、路径最短，白天很短。
\end{solution}

\begin{solution}{exo:g3:sun-days-seasons:2}
正午时太阳最高。正午的影子也
每天、每年都指向同一方向——
是大自然中最可靠的“箭头”。
\end{solution}

\begin{solution}{exo:g3:sun-days-seasons:3}
六月，太阳走高路，正午时几乎在头顶，
所以影子短短的，像一截树桩。十二月，太阳走低路，
即使在正午，阳光也是斜射的，
同一根柱子投下长长的影子。
\end{solution}

\begin{solution}{exo:g3:sun-days-seasons:4}
太阳沿弧线移动时，杆子的影子会扫过
各个刻度；看影子碰到哪条刻度线，就能读出时间。
移动的“指针”就是影子本身。
\end{solution}

\begin{solution}{exo:g3:sun-days-seasons:5}
若某个时刻墙的阴影把杆子整个吞没，
那个时刻就无法标记和读数。阴天时，日晷没有“指针”——
没有影子；夜里也没有：日晷只在
太阳照耀时才能使用。
\end{solution}

\begin{solution}{exo:g3:sun-days-seasons:6}
一年约 $365$ 天，绕行一圈有四个季节：
春、夏、秋、冬。
\end{solution}

\begin{solution}{exo:g3:sun-days-seasons:7}
一年约 $365$ 天；一天有 $24$ 小时。
一天是地球绕自身转一周，即\emph{自转}；一年是地球绕太阳转一周，即\emph{公转}。
这两种运动各不相同。
\end{solution}

\begin{solution}{exo:g3:sun-days-seasons:8}
例如：路灯亮得多早（晚饭前就天黑，说明是冬天）；
正午太阳有多高；正午的
影子有多长；晚饭后能否在天亮时到户外玩耍。
\end{solution}

\begin{solution}{exo:g3:sun-days-seasons:9}
同一时刻他们正值夏天。若整个地球
都离太阳更远，则\emph{所有人}会同时进入冬天——
所以这个猜想必定不对。真正的原因是太阳走低路，
以及它斜射的光线。
\end{solution}

\begin{solution}{exo:g3:sun-days-seasons:10}
杆子没有动——动的是太阳的路线。从三月到六月，
太阳的弧线越爬越高，影子变短并且移动了，所以
现在的影子已经碰不到为三月路线铺设的小石子。日晷
必须结合季节来读数。
\end{solution}

\begin{solution}{exo:g3:sun-days-seasons:11}
太阳最高时，说明她是正午醒来的。若再过一小时
帐篷的影子就已经比帐篷还长，说明太阳
即使接近正午也挂得很低——这是冬季的低路：
她所在的日子更接近隆冬。
\end{solution}
